\section{5. Results}\label{results}

\subsection{5.1 Nominal fitted-model
dependence}\label{nominal-fitted-model-dependence}

The nominal fitted-model results are summarized below.

\begin{longtable}[]{@{}
  >{\raggedright\arraybackslash}p{(\columnwidth - 6\tabcolsep) * \real{0.2000}}
  >{\raggedleft\arraybackslash}p{(\columnwidth - 6\tabcolsep) * \real{0.2667}}
  >{\raggedleft\arraybackslash}p{(\columnwidth - 6\tabcolsep) * \real{0.2667}}
  >{\raggedleft\arraybackslash}p{(\columnwidth - 6\tabcolsep) * \real{0.2667}}@{}}
\toprule\noalign{}
\begin{minipage}[b]{\linewidth}\raggedright
Specimen
\end{minipage} & \begin{minipage}[b]{\linewidth}\raggedleft
Boundary-excluded span (meV)
\end{minipage} & \begin{minipage}[b]{\linewidth}\raggedleft
Full completeness span (meV)
\end{minipage} & \begin{minipage}[b]{\linewidth}\raggedleft
Maximum coherent calibration-corner shift (meV)
\end{minipage} \\
\midrule\noalign{}
\endhead
\bottomrule\noalign{}
\endlastfoot
MZ226 & 5.09 & 6.41 & 0.89 \\
MZ310 & 6.68 & 6.83 & 0.68 \\
\end{longtable}

The primary result is the boundary-excluded span. The larger
completeness span includes fits pinned at an edge-search bound and is
not interpreted as an identified physical range.

\begin{figure}
\centering
\includegraphics{figures/figure2_spectrum_operator_overlays.pdf}
\caption{Reconstructed MZ226 spectrum and fitted extraction models.
Fitted curves are shown only over the declared nominal fit population.}
\end{figure}

\subsection{5.2 Local irregularity and contiguous-window
leverage}\label{local-irregularity-and-contiguous-window-leverage}

Omitting the two MZ226 coordinates associated with the visible plateau
changes a non-boundary fitted intercept by at most approximately 0.12
meV. Combining the omission with the calibration-corner cases produces a
maximum change of approximately 0.96 meV. The remaining non-boundary
model span is approximately 5.12 meV. The local pair therefore does not
account for the primary model-form dependence.

The broader leverage test is more consequential.

\begin{longtable}[]{@{}
  >{\raggedright\arraybackslash}p{(\columnwidth - 8\tabcolsep) * \real{0.1579}}
  >{\raggedleft\arraybackslash}p{(\columnwidth - 8\tabcolsep) * \real{0.2105}}
  >{\raggedleft\arraybackslash}p{(\columnwidth - 8\tabcolsep) * \real{0.2105}}
  >{\raggedleft\arraybackslash}p{(\columnwidth - 8\tabcolsep) * \real{0.2105}}
  >{\raggedleft\arraybackslash}p{(\columnwidth - 8\tabcolsep) * \real{0.2105}}@{}}
\toprule\noalign{}
\begin{minipage}[b]{\linewidth}\raggedright
Specimen
\end{minipage} & \begin{minipage}[b]{\linewidth}\raggedleft
Leave-one-out max (meV)
\end{minipage} & \begin{minipage}[b]{\linewidth}\raggedleft
Three-point max (meV)
\end{minipage} & \begin{minipage}[b]{\linewidth}\raggedleft
Five-point max (meV)
\end{minipage} & \begin{minipage}[b]{\linewidth}\raggedleft
Minimum non-boundary span across tested omissions (meV)
\end{minipage} \\
\midrule\noalign{}
\endhead
\bottomrule\noalign{}
\endlastfoot
MZ226 & 0.90 & 2.33 & 3.60 & 4.98 \\
MZ310 & 1.03 & 3.01 & 4.51 & 3.39 \\
\end{longtable}

The largest shifts occur in the free-exponent model when the earliest
admitted points are removed. This is evidence of practical
lower-boundary conditioning, not evidence that the corresponding
coordinates are incorrect.

\subsection{5.3 Reconstruction, fit-domain, weighting, and residual
compatibility}\label{reconstruction-fit-domain-weighting-and-residual-compatibility}

The committed unweighted isotonic reconstruction is reproducible from
the stored source-pixel centers to within 0.0085 and 0.0032 in
\(\log_{10}\alpha\) for MZ226 and MZ310, respectively. With nominal fit
membership held fixed, the raw centerline, independently recomputed
isotonic curve, inverse-width-weighted isotonic curve, and lower/upper
line-envelope reconstructions move any fitted intercept by at most 0.69
meV for MZ226 and 0.81 meV for MZ310. When each reconstruction is
allowed to determine its own fit membership, the largest admissible
shifts are 1.25 and 1.96 meV.

One upper-line-envelope membership scenario for each specimen is
excluded rather than repaired: its lowest admitted photon energy falls
below the predeclared upper edge-search bound. Changing the search bound
to retain those scenarios would introduce an unplanned analyst choice.

The fit-domain and weighting grid is substantially more influential.
Twenty-seven admissible combinations are evaluated for MZ226 and 24 for
MZ310. Across them, the largest displacement of an individual fitted
intercept is 5.71 meV for MZ226 and 5.44 meV for MZ310, both in the
free-exponent model. The non-boundary complete-model span ranges from
2.59 to 11.00 meV for MZ226 and from 1.84 to 11.09 meV for MZ310. Thus
no single nominal model-span value should be treated as invariant to the
declared fit convention.

\begin{longtable}[]{@{}
  >{\raggedright\arraybackslash}p{(\columnwidth - 6\tabcolsep) * \real{0.2000}}
  >{\raggedleft\arraybackslash}p{(\columnwidth - 6\tabcolsep) * \real{0.2667}}
  >{\raggedleft\arraybackslash}p{(\columnwidth - 6\tabcolsep) * \real{0.2667}}
  >{\raggedleft\arraybackslash}p{(\columnwidth - 6\tabcolsep) * \real{0.2667}}@{}}
\toprule\noalign{}
\begin{minipage}[b]{\linewidth}\raggedright
Specimen
\end{minipage} & \begin{minipage}[b]{\linewidth}\raggedleft
Reconstruction, fixed membership (meV)
\end{minipage} & \begin{minipage}[b]{\linewidth}\raggedleft
Reconstruction + membership (meV)
\end{minipage} & \begin{minipage}[b]{\linewidth}\raggedleft
Endpoint/weighting max (meV)
\end{minipage} \\
\midrule\noalign{}
\endhead
\bottomrule\noalign{}
\endlastfoot
MZ226 & 0.69 & 1.25 & 5.71 \\
MZ310 & 0.81 & 1.96 & 5.44 \\
\end{longtable}

\begin{longtable}[]{@{}
  >{\raggedright\arraybackslash}p{(\columnwidth - 2\tabcolsep) * \real{0.4286}}
  >{\raggedleft\arraybackslash}p{(\columnwidth - 2\tabcolsep) * \real{0.5714}}@{}}
\toprule\noalign{}
\begin{minipage}[b]{\linewidth}\raggedright
Specimen
\end{minipage} & \begin{minipage}[b]{\linewidth}\raggedleft
Non-boundary fitted-model span across endpoint/weighting grid (meV)
\end{minipage} \\
\midrule\noalign{}
\endhead
\bottomrule\noalign{}
\endlastfoot
MZ226 & 2.59-11.00 \\
MZ310 & 1.84-11.09 \\
\end{longtable}

The heatmaps below report the largest absolute free-exponent
displacement across the three weighting rules at each endpoint pair.
Bracketed values give the signed range across those weighting rules. The
response varies systematically with the lower and upper endpoints rather
than being caused by a single anomalous grid cell.

\begin{figure}
\centering
\includegraphics{figures/figure8_fit_grid_mz226.pdf}
\caption{MZ226 free-exponent fit-domain and weighting sensitivity.}
\end{figure}

\begin{figure}
\centering
\includegraphics{figures/figure9_fit_grid_mz310.pdf}
\caption{MZ310 free-exponent fit-domain and weighting sensitivity.}
\end{figure}

\FloatBarrier

\begin{figure}
\centering
\includegraphics{figures/figure4_robustness_scales.pdf}
\caption{Maximum fitted-intercept displacement under distinct
deterministic robustness analyses.}
\end{figure}
